\documentclass[11pt]{article}
\textwidth 15.9cm
\textheight 23. cm
\addtolength{\oddsidemargin}{-11.5mm}
\addtolength{\evensidemargin}{-11.5mm}
\addtolength{\topmargin}{-19.0mm}

%\usepackage{amsmath}
\usepackage[ansinew]{inputenc}
\usepackage[bbgreekl]{mathbbol}
\usepackage{amsmath,amssymb}
\usepackage{graphicx}
\usepackage{hyperref}
\usepackage{url}
\usepackage[numbers,square]{natbib}
\usepackage{multirow}
\usepackage{tikz}
\usepackage{tikz-cd}
\usetikzlibrary{arrows,matrix}
\usepackage{bbm}
\usepackage{authblk}
\usepackage{subcaption}
\renewcommand{\Affilfont}{\footnotesize}
%\renewcommand{\Authfont}{\normalsize}

\usepackage{stmaryrd}
\usepackage{comment}
\usepackage{showlabels}
\usepackage{soul}
%\usepackage[dvipsnames]{xcolor}

\usetikzlibrary{calc,patterns,shapes.geometric}
\usetikzlibrary{arrows,shapes,positioning}
\usetikzlibrary{decorations.markings,decorations.pathmorphing,decorations.pathreplacing}
\usetikzlibrary{decorations.text,decorations.shapes}

%%%%%%%% MACROS
\include{gempic_spectral_macros}

\title{Exterior Derivatives}


\author[1]{The Gempic development team}
\affil[1]{Max-Planck-Institut f\"ur Plasmaphysik, Garching, Germany}


\begin{document}

\maketitle

\begin{abstract}
This note describes some conventions that should be followed in the GEMPIC code
and its documentation.
\end{abstract}

\tableofcontents

\section{Curl}
The curl is calculated applying the Stokes' theorem along all discrete elements of the domain. For example, if $\bfE$ is a 3D-vector defined in $\RR^3$, $\bff^1_{I} = (f^1_{I,x}, f^1_{I,y}, f^1_{I,z})$ the projections of $\bff$ to 1-forms (line integrals along each direction) and $\bff^2_{I} = (f^2_{I,x}, f^2_{I,y}, f^2_{I,z})$ the corresponding projections of $\bff$ to 2-forms (face integrals on the elements of the grid), then:

\begin{equation}
    \oint_{\partial \Sigma} \bff \cdot d\bfl = \iint_{\Sigma} \curl f \cdot d\bfS
\end{equation}

Which means that we can calculate the curl with the following relations applied to each component of the 2-form:

\begin{align}
    f^2_{I,x}(x, y, z) = f^1_{I,z}(x, y + dy, z) - f^1_{I,z}(x, y, z) - f^1_{I,y}(x, y, z + dz) + f^1_{I,y}(x, y, z)
\end{align}

\begin{align}
    f^2_{I,y}(x, y, z) = f^1_{I,x}(x, y, z + dz) - f^1_{I,x}(x, y, z) - f^1_{I,z}(x + dx, y, z) + f^1_{I,z}(x, y, z)
\end{align}

\begin{align}
    f^2_{I,z}(x, y, z) = f^1_{I,y}(x + dx, y, z) - f^1_{I,y}(x, y, z) - f^1_{I,x}(x, y + dy, z) + f^1_{I,z}(x, y, z)
\end{align}

\section{Div}
To be included.

\section{Grad}
To be included.

\end{document}
